% DRAFT skeleton, started 2026-10-24.  IEEE conference format (IEEEtran).  NOT compiled: no TeX distribution is installed on the development PC (see README.md).
% Rules for this file: every number comes from a stored result file named in a comment next to it; anything not measured on the board is marked "pending board run".
% Baseline everywhere = the DEPLOYED checkpoint (INT8 PSNR-Y 35.02 / 31.40 / 30.53 dB), not the 2-seed R0 mean of the ladder study.
% Venue and deadline are not chosen; page limit and review type (single / double blind) are not known.  A public repository link must be anonymised for double-blind review.
\documentclass[conference]{IEEEtran}
\usepackage{cite}
\usepackage{amsmath}
\usepackage{graphicx}
\usepackage{booktabs}
\usepackage{url}

\begin{document}

\title{An INT8 Super-Resolution Accelerator on a Low-Cost Zynq-7020:\\Bit-Exact Verification and a Measured Utilization Gap}

\author{\IEEEauthorblockN{TODO: authors}
\IEEEauthorblockA{TODO: affiliation}}

\maketitle

\begin{abstract}
% Only measured statements; edit after the board session.
We describe a $\times 2$ image super-resolution accelerator for a ZedBoard (XC7Z020): a 2{,}620-parameter INT8 CNN upscales $960\times540$ to $1920\times1080$
in 1.312\,s per frame (0.762\,fps, 20-frame mean, board-measured) with 68 DSP48 blocks at 100\,MHz.
The output is byte-identical to a NumPy integer model on all 119 images of Set5, Set14 and BSD100 on the board.
On these sets the INT8 network gains about 1.4, 1.1 and 1.0\,dB PSNR-Y over bicubic interpolation (Table~\ref{tab:quality}).
TODO: contributions after the board session (utilization gap and its closure, mutation-tested RTL, energy against the on-chip ARM core).
\end{abstract}

\begin{IEEEkeywords}
super-resolution, FPGA, Zynq, INT8 quantization, bit-exact verification
\end{IEEEkeywords}

\section{Introduction}
% TODO: motivation (low-cost edge FPGA, why a tiny network), problem statement, contributions (to be fixed AFTER the related-work reading; see docs/paper/related_work_and_venues.md:
% the layer pipeline and bit-exact golden-model verification are NOT new ideas).
TODO.

\section{Related Work}
% TODO: read each paper in full before citing (docs/paper/related_work_and_venues.md section 1 lists the candidates: FPGA SR systems for UHD video, real-time SR on edge FPGAs,
% energy-efficient deconvolution accelerators, CGRA4ML for bit parity, mutation testing of HDL).  Only FSRCNN has been read so far.
Compact networks such as FSRCNN~\cite{dong2016fsrcnn} reach good quality with few parameters; our network follows the same idea at a smaller scale.
TODO: FPGA super-resolution accelerators; layer-pipelined designs; verification against golden models.

\section{Model}\label{sec:model}
The network (\texttt{software/ai/models/srnet.py}) maps an RGB image to a $2\times$ larger one with four convolutions and a pixel shuffle:
$3\times3$ convolution $3\!\to\!16$ with ReLU, $3\times3$ depthwise convolution (16 channels) with ReLU, $1\times1$ convolution $16\!\to\!16$ with ReLU,
$3\times3$ convolution $16\!\to\!12$, and a pixel shuffle with factor 2 (channel $k=4c+2\,dy+dx$).
All convolutions are valid (no padding): the input carries a 3-pixel border (replicated for whole images, real neighbours for tiles).
It has 2{,}620 parameters and 2{,}560 multiply-accumulates (MACs) per input pixel (asserted in \texttt{software/ai/evaluation/test\_model.py}).

\begin{table}[t]
\caption{Layers of the network. Parameters include biases.}\label{tab:layers}
\centering
\begin{tabular}{llrr}
\toprule
Layer & Shape & Parameters & MACs / input pixel \\
\midrule
L1 conv $3{\times}3$ + ReLU & $3\to16$ & 448 & 432 \\
L2 depthwise $3{\times}3$ + ReLU & 16 & 160 & 144 \\
L3 conv $1{\times}1$ + ReLU & $16\to16$ & 272 & 256 \\
L4 conv $3{\times}3$ & $16\to12$ & 1{,}740 & 1{,}728 \\
\midrule
Total & & 2{,}620 & 2{,}560 \\
\bottomrule
\end{tabular}
% per-layer values computed from the layer shapes: L1 3*3*3*16 (+16), L2 16*9 (+16), L3 16*16 (+16), L4 16*9*12 (+12); totals equal the asserted 2,620 / 2,560.
\end{table}

Training uses 180 DIV2K images with an L1 loss (the checkpoint \texttt{srnet\_fp32.pt}; \texttt{results/quality/m2\_training\_summary.md}).
The low-resolution inputs are made by Pillow bicubic downscaling after cropping the high-resolution image to even sides; PSNR and SSIM are computed on the luma channel (BT.601)
with a 2-pixel border shaved, SSIM with a Gaussian $11\times11$ window, $\sigma=1.5$ (\texttt{software/ai/evaluation/metrics.py}).
TODO: training details to be checked against \texttt{train.py} before submission.

\section{Quantization}\label{sec:quant}
Activations are unsigned 8-bit with one scale per layer, weights are signed 8-bit with a per-layer scale (values in $[-127,127]$), biases and accumulators are 32-bit.
A layer output is $y=\mathrm{clamp}\big((acc\cdot M+2^{S-1})\gg S,\,0,\,255\big)$ with a 16-bit multiplier $M$ per output channel and a shift $S$ per layer
($S=24,23,21,24$ for L1--L4); the product $acc\cdot M$ needs about 40 bits.
The activation range of each layer is a percentile of all post-ReLU values over 100 training crops.
The percentile was chosen on 12 validation images (Table~\ref{tab:ptq}); post-training quantization (PTQ) is sufficient, quantization-aware training was not needed for the deployed network.

\begin{table}[t]
\caption{PTQ percentile sweep, per-layer weight scales, 12 validation images (FP32 reference 34.954\,dB PSNR-Y).}\label{tab:ptq}
\centering
\begin{tabular}{rrr}
\toprule
Percentile & PSNR-Y INT8 (dB) & Drop vs FP32 (dB) \\
\midrule
99.0 & 29.648 & 5.307 \\
99.9 & 34.529 & 0.425 \\
\textbf{99.99} & \textbf{34.731} & \textbf{0.223} \\
100.0 & 34.675 & 0.280 \\
\bottomrule
\end{tabular}
% source: results/quality/m3_ptq_sweep.md (per-channel weights were slightly worse in the same file and were not used)
\end{table}

On the test sets the INT8 model loses on average 0.182, 0.111 and 0.076\,dB PSNR-Y against the FP32 network on Set5, Set14 and BSD100
(worst single image 0.314, 0.285 and 0.413\,dB; \texttt{results/quality/model\_int8\_x2.md}).
The INT8 model is implemented as a NumPy integer model (\texttt{integer\_reference.py}); it is the golden reference of the hardware (Section~\ref{sec:arch}).

\section{Architecture}\label{sec:arch}
% TODO: figure (PS, DMA, AXI-Lite registers, tile core with four engines), text.  Facts to use (all stored):
%  - 64x64 core tile + 3 px halo = 70x70 input (14,700 B in, 49,152 B out per tile), 960x540 padded to 960x576 = 135 tiles          [HANDOFF section 3]
%  - one convolution engine per layer, run one after another; 926,201 cycles per tile in simulation (engines 854,149), 70 BRAM tiles and 68 DSP for the core [local project log 6e/6f, results/utilization/m5_tile_core.md]
%  - 32-bit AXI-Stream wrapper with TLAST and error flags, AXI4-Lite registers, ARM driver with timeouts and error recovery        [docs/REGISTER_MAP.md]
TODO.

\section{Results}\label{sec:results}
All board numbers were measured on a ZedBoard with the deployed bitstream (VERSION 0x53520602, 100\,MHz), the test image is a DIV2K image of $960\times540$.
Power is a Vivado estimate, not a measurement.

\begin{table}[t]
\caption{Image quality, $\times2$, PSNR-Y (dB) / SSIM-Y. FPGA = outputs read back from the board, byte-identical to the INT8 integer model.}\label{tab:quality}
\centering
\begin{tabular}{lccc}
\toprule
Method & Set5 & Set14 & BSD100 \\
\midrule
Bicubic & 33.67 / 0.9303 & 30.32 / 0.8698 & 29.56 / 0.8434 \\
FP32 network & 35.21 / 0.9449 & 31.51 / 0.8975 & 30.60 / 0.8785 \\
INT8 integer model & 35.02 / 0.9420 & 31.40 / 0.8948 & 30.53 / 0.8756 \\
FPGA (board) & 35.02 / 0.9420 & 31.40 / 0.8948 & 30.53 / 0.8756 \\
\bottomrule
\end{tabular}
% source: results/quality/table_a_x2.md (119 of 119 images byte-exact on the board; INT8 beats bicubic on 119 of 119 images)
\end{table}

Table~\ref{tab:quality} shows the quality; the network is better than bicubic interpolation on every one of the 119 images, and the board output equals the integer model on every image.
Fine random texture (gravel, fur, fabric) is not recovered by a network of this size.

\begin{table}[t]
\caption{Candidate weight sets against the deployed design (R0). Quality: INT8 integer model, PSNR-Y (dB), the single built seed of each candidate. Resources and timing: Vivado post-implementation of the whole system at 100\,MHz, not board measurements.}\label{tab:candidates}
\centering
\begin{tabular}{lcccc}
\toprule
Design & Set5 / Set14 / BSD100 & LUT & WNS (ns) & Board \\
\midrule
R0 deployed & 35.02 / 31.40 / 30.53 & 5{,}296 & +0.503 & 0.762\,fps \\
R2q (rep.\ + QAT) & 35.84 / 31.89 / 30.87 & 5{,}287 & +0.944 & pending board run \\
R3q (anchor + QAT) & 36.20 / 32.12 / 31.01 & 5{,}323 & +0.683 & pending board run \\
\bottomrule
\end{tabular}
% sources: R0 results/quality/table_a_x2.md, results/utilization/m6_system_0602.md, results/speed/m7_timing_20frames.md;
%          R2q / R3q: results/utilization/m6_system_R2q.md (branch r2q), m6_system_R3q.md (branch r3q) -- INT8 lines and the build tables.
% R2q needs only different constants (SHIFT 22/22/21/22); R3q adds a nearest-neighbour anchor (+33..45 LUT, +66 FF, no extra cycle: m6_system_R3q.md).
\end{table}

\begin{table}[t]
\caption{Speed, resources and power of the deployed design (R0), $960\times540\to1920\times1080$.}\label{tab:system}
\centering
\begin{tabular}{ll}
\toprule
Item & Value \\
\midrule
Time per frame (20-frame mean) & 1{,}311.9\,ms (min 1{,}311.4, max 1{,}313.2) \\
Frames per second & 0.762 \\
Cycles per 64$\times$64 tile (board counter) & 926{,}217--926{,}241 (simulation 926{,}202) \\
LUT / FF / BRAM tiles / DSP & 5{,}296 / 7{,}831 / 72 / 68 \\
Timing at 100\,MHz (WNS / WHS) & +0.503 / +0.016\,ns \\
Power, Vivado estimate (chip / accelerator) & 1.913 / 0.176\,W, confidence low \\
ARM-only baseline (same chip, no accelerator) & pending board run \\
\bottomrule
\end{tabular}
% sources: results/speed/m7_timing_20frames.md; results/utilization/m6_system_0602.md; results/utilization/m7_power.md.
% 95.3 % of the frame is the core itself (1,250.4 ms of 1,311.9 ms, m7_timing_20frames.md).
\end{table}

\begin{table}[t]
\caption{The same network on a desktop PC (not a fair energy comparison; PC times are medians over repeated runs).}\label{tab:pc}
\centering
\begin{tabular}{lrr}
\toprule
Implementation & ms / frame & fps \\
\midrule
FP32 PyTorch, CPU, 1 thread & 128.8 & 7.76 \\
FP32 PyTorch, GPU incl.\ host copies & 7.9 & 126.81 \\
INT8 NumPy integer model, CPU (unoptimised) & 1{,}526.2 & 0.66 \\
ZedBoard FPGA (board-measured) & 1{,}311.9 & 0.76 \\
\bottomrule
\end{tabular}
% source: docs/paper/cpu_baseline.md (i5-13450HX, RTX 3050 6 GB laptop GPU, PyTorch 2.12.0+cu130; 2026-10-16).  CPU energy not measured; GPU 0.548 J per frame is GPU power only.
\end{table}

\section{Discussion}
% TODO after the board session and the fmax decision.  Facts to discuss (all stored):
%  - utilization gap: only one of the four engines works at a time, 17.2 % of the 68 DSPs busy (cycle MODEL, docs/paper/cycle_model.md Part 1); layer 4 is 64 % of the cycles.
%  - predictions (not built): line-buffer layer pipeline about 3.9 fps at 68 DSP, 100 MHz (Part 2); shared engine 2.2 / 2.8 fps (Part 3).  Label as predictions.
%  - 150 MHz does not close timing without a pipeline register (docs/paper/timing_150mhz.md; reproduction status there).
%  - limits: 0.76 fps is not real time; load, compute and output are not overlapped; power is an estimate; images go in over JTAG; two single first-tile failures of 2026-10-11 remain unexplained.
TODO.

\section*{Acknowledgment}
TODO.

\bibliographystyle{IEEEtran}
\bibliography{refs}

\end{document}
